\section{H54: The kickoff text sets the day-1 target}

\textbf{The question:} When a new goal period starts, the operators post a kickoff text. Does the swarm's first-day talk land on what that text names? And does a more specific text give a tighter swarm?

\textbf{The intuition:} Think of a quench: a sudden change of a control parameter, after which a system relaxes toward a new state. Each agent's statements form a vector (a point in a meaning space). The kickoff acts as a field, a common pull on every agent. H54 says the kickoff text itself gives the direction of that pull. Each agent also keeps its own offset (style, model family, old project) and feels some herding from the others. So the day-1 average of all statements should sit closer to its own kickoff than to any other kickoff. In the Potts picture, each agent chooses one of several discrete projects. There the kickoff is a field on the projects it names, and those projects freeze at once.

\textbf{What we measured:} For each of 33 non-reserved goal periods, we took the day-1 content centroid (the mean statement vector of the first day). We scored it against its own kickoff and the 32 other kickoffs. The null is a kickoff swap: if the text sets no target, the own kickoff ranks at random. A genericness correction removes credit from kickoffs that resemble every centroid. Stricter decoys were kickoffs from the same regime and from neighbouring periods. The previous period's last-day state tested inertia. The scheduler is not relevant, because the test uses day-level content and no timing. The external drive is the object of study. Style-residualized vectors removed the shared model prior (top-1 stays at 0.52). Common input (agents quoting the kickoff) is part of the mechanism; removing near-echo statements tests how much. A synthetic check came first: vector-spin swarms at the real day-1 counts. Under no target the test passed 0 of 200 times; with a moderately faithful text proxy it recovered the target.

\textbf{What we found:} The day-1 centroid picks out its own kickoff. The median own-kickoff percentile is 1.0, and the own kickoff ranks first in 18 of 33 periods ($p=3\times10^{-6}$). The day-1 move from the previous period points at the new kickoff in 23 of 25 periods. The test fails only where the kickoff names no shared target: free weeks \#3 and \#16, \#44's half-free week and \#51. In \#51 each agent has a private goal text, and each agent lands on its own goal (role-swap accuracy 0.95, $p=0.0002$, flat over nine weeks). A human reassigned Claude Opus~5 on 07-29 (NE38). The agent moved onto its new goal within a day (difference-in-differences $+0.61$ [0.56, 0.66]). Nine of 13 kickoff-frozen projects carry the goal text's words (enrichment 3.2, stratified $p=0.006$); only one pre-existed. Text specificity does not set the spread ($\rho=+0.17$; the prediction was $\le-0.35$). Scope: 33 kickoffs (22 in regime~I), plus \#51 and NE38. The claim that stands is ``day-1 content lands on what the kickoff names; agents land on private goals''\cred{0.92}, EU 2.6. The summary metadata and the card agree.

\begin{figure}[h]
\includegraphics[width=\textwidth]{H54-kickoff-quench-target/fig.pdf}
\caption{Day-1 content lands on its own kickoff text. (a) Each row is one goal period's day-1 centroid, scored against all 33 kickoffs; the diagonal is the own kickoff. (b) The own kickoff ranks first in 18 of 33 periods. The gray band is the swap null; the misses are kickoffs with no shared target. (c) In \#51 each agent sits closer to its own private goal than to another agent's goal in 0.95 of pairs, for nine weeks.}
\end{figure}

\textbf{What it means:} For an operator: write the shared target into the kickoff text. Day-1 talk ranks its own kickoff first in 18 of 33 periods and in the top tenth in 26 of 33. The frozen project usually carries the goal's words (9/13). To move one agent, rewrite that agent's own goal text: in NE38 the agent moved within one day. Do not expect a ``more specific'' prompt (more numbers, names or deadlines) to tighten the swarm; the measured link is zero. A choose-your-own-goal kickoff gives no shared target and a more scattered swarm. Agent-written plans, including an elected leader's announcement, did not act as targets.

\textbf{Caveats:} The target is a text-embedding proxy from one model (bge-small); the second embedding model was not run. Quoting is part of the effect: removing near-echo statements lowers top-1 from 0.55 to 0.45 or 0.30, but the median percentile stays $\ge0.9$. ``Named'' uses name-token matching, and agents may name repositories after the goal, so the direction of naming is open. The frozen-project test (P2) is mixed by its literal rule and underpowered. The two-room tests in \#44 and \#38 failed and were underpowered. The free-week explanation is post hoc. The confirmation script is written and dry-run; no confirmation test on reserved goal periods has run yet.
